\documentclass{article} 
\usepackage{graphicx}
\usepackage[francais]{babel} 
\usepackage[utf8]{inputenc}
\usepackage[T1]{fontenc} 
\usepackage{amsmath} 
\usepackage{amsfonts}
\usepackage{verbatim} 
\usepackage{float} 
\usepackage{hyperref}
\usepackage{scrtime}
\usepackage[scale=2]{ccicons}
\usepackage{tabularx}
\usepackage{Sweave}
\begin{document}
\Sconcordance{concordance:serie2-solutions.tex:/home/francois/ACT-2010-Exercices/serie2-solutions.Rnw:%
1 13 1 1 0 130 1}

%%% Police de caractère pour la page titre
\renewcommand{\sfdefault}{hls}

%%% Marge de gauche (1/3 de la page)
\newlength{\gauche}
\addtolength{\gauche}{72mm}
\addtolength{\gauche}{-\spinemargin}

%%% Épaisseur de la bande sur la page titre
\newlength{\ruleheight}
\setlength{\ruleheight}{7.75mm}

%%% Définition de la bande
\definecolor{gray}{gray}{0.4}
\textblockorigin{0mm}{279mm}
\newcommand{\banderecto}{%
  \begin{textblock*}{71.5mm}[0,1](0mm,-46.5mm)
    \textblockcolor{black} \rule{0mm}{\ruleheight}
  \end{textblock*}
  \begin{textblock*}{144mm}[0,1](72mm,-46.5mm)
    \textblockcolor{gray} \rule{0mm}{\ruleheight}
  \end{textblock*}}
\newcommand{\bandeverso}{%
  \begin{textblock*}{144mm}[0,1](0mm,-46.5mm)
    \textblockcolor{gray} \rule{0mm}{\ruleheight}
  \end{textblock*}
  \begin{textblock*}{71.5mm}[0,1](144.5mm,-46.5mm)
    \textblockcolor{black} \rule{0mm}{\ruleheight}
  \end{textblock*}}

%%% Titre
\begin{adjustwidth*}{\gauche}{-15mm}
  \sffamily\fontseries{ub}\selectfont
  \raggedright
  \vspace*{2cm}
  \thetitle
\end{adjustwidth*}

%%% Affichage de la bande
\banderecto
\cleardoublepage

%%% Page de garde
\begin{adjustwidth*}{\gauche}{-15mm}
  \sffamily\fontseries{ub}\selectfont
  \raggedright
  \vspace*{2cm}
  \thetitle \\
  \bfseries
  \vspace*{3cm}
  \theauthor
  \vspace*{\fill}
  \thedate
\end{adjustwidth*}

\clearpage

%%% Page de notices
\begingroup
\calccentering{\unitlength}
\begin{adjustwidth*}{\unitlength}{-\unitlength}
  \small
  \setlength{\parindent}{0pt}
  \setlength{\parskip}{\baselineskip}

  {\textcopyright} 2013 François Pelletier \\


  \includegraphics[height=7mm,keepaspectratio=true]{by}\;%
  \includegraphics[height=7mm,keepaspectratio=true]{sa} \\
  Cette création est mise à disposition selon le contrat
  Paternité-Partage des conditions initiales à l'identique 2.5 Canada
  disponible en ligne
  \url{http://creativecommons.org/licenses/by-sa/2.5/ca/} ou par
  courrier postal à Creative Commons, 171 Second Street, Suite 300,
  San Francisco, California 94105, USA.

  \textbf{Historique de publication}
  \vspace{-\baselineskip}
  \begin{tabbing}
    Décembre 2013:\quad\= Première édition
  \end{tabbing}

  \textbf{Code source} \\
  Le code source {\LaTeX} de ce document est disponible à l'adresse
  \begin{quote}
    \url{https://github.com/franc00018/ACT-2010-Exercices}
  \end{quote}
  ou en communiquant directement avec les auteurs.

  \vspace{1cm}

  % ISBN \ISBN \\
  % Dépôt légal -- Bibliothèque et Archives nationales du Québec, 2014 \\
  % Dépôt légal -- Bibliothèque et Archives Canada, 2014
\end{adjustwidth*}
\endgroup

%%% Retour à la police normale
\renewcommand{\sfdefault}{phv}

\clearpage

%%% Local Variables:
%%% mode: latex
%%% TeX-master: "exercices_series_chrono"
%%% coding: utf-8-unix
%%% End:

\addtocounter{section}{1}
\section{Modèles classiques pour séries chronologiques}

\subsection{Zone de stationnarité pour AR(2) (Théorique)}

On obtient les racines de l'équation caractéristique en utilisant la formule quadratique habituelle

\begin{align*}
  \frac{\phi_1\pm\sqrt{\phi_1^2+4\phi_2}}{-2\phi_2}
\end{align*}

On considère les inverses des deux racines, $A_1$ et $A_2$.

\begin{align*}
  A_1 &= \frac{2\phi_2}{-\phi_1-\sqrt{\phi_1^2+4\phi_2}} \\
  &= \frac{2\phi_2}{-\phi_1-\sqrt{\phi_1^2+4\phi_2}} \left[\frac{-\phi_1+\sqrt{\phi_1^2+4\phi_2}}{-\phi_1+\sqrt{\phi_1^2+4\phi_2}} \right] \\
  &= \frac{2\phi_2(-\phi_1+\sqrt{\phi_1^2+4\phi_2})}{\phi_1^2-(\phi_1^2+4\phi_2)}\\
  &= \frac{\phi_1-\sqrt{\phi_1^2+4\phi_2}}{2}\\
  A_2 &= \frac{\phi_1+\sqrt{\phi_1^2+4\phi_2}}{2}
\end{align*}

Il y a 2 situations possibles: soit les racines sont réelles ($\phi_1^2+4\phi_2>0$) ou elles sont complexes ($\phi_1^2+4\phi_2<0$).
\begin{itemize}
\item \textbf{Racines réelles:}
  Comme les racines doivent être plus grandes que 1, alors nécessairement leurs inverses $|A_1|<1$ et $|A_2|<1$. Nous avons donc:
  \begin{align*}
    -1 &< \frac{\phi_1-\sqrt{\phi_1^2+4\phi_2}}{2} <  \frac{\phi_1+\sqrt{\phi_1^2+4\phi_2}}{2} < 1 \\
    \Leftrightarrow -2 &< \phi_1-\sqrt{\phi_1^2+4\phi_2} < \phi_1+\sqrt{\phi_1^2+4\phi_2} < 2
  \end{align*}
  En observant la première inégalité, on a:
  \begin{align*}
    -2 < \phi_1-\sqrt{\phi_1^2+4\phi_2}
    &\Leftrightarrow \sqrt{\phi_1^2+4\phi_2}<\phi_1+2 \\
    &\Leftrightarrow \phi_1^2+4\phi_2 < \phi_1^2+4\phi_1+4 \\
    &\Leftrightarrow \phi_2 < \phi_1 + 1 \\
    &\Leftrightarrow \phi_2 - \phi_1 < 1
  \end{align*}
  Ce qui correspond à la seconde condition. En considérant la seconde inégalité, on obtient, de la même façon, la première inégalité:
  \begin{align*}
    \phi_1+\sqrt{\phi_1^2+4\phi_2} &< 2 \\
    &\Leftrightarrow \phi_1 + \phi_2 < 1
  \end{align*}
  Ces deux conditions réunies avec un discriminant positif forment la région de stationnarité pour des racines réelles.
  
\item \textbf{Racines complexes:}  
  On considère la situation où $\phi_1^2+4\phi_2<0$. Ici, on aura des conjugués complexes et $|A_1| = |A_2| <1$ seulement si $|A_1|^2<1$.
  \begin{align*}
    |A_1|^2 &=\frac{\phi_1^2+(-\phi_1^2-4\phi_2}{4}=-\phi^2 \\
    &\Leftrightarrow \phi_2>-1 \\
    &\Leftrightarrow |\phi_2|<1
  \end{align*}
  Ce résultat réuni avec un discriminant négatif forment la région de stationnarité pour des racines complexes.
\end{itemize}

\clearpage
\subsection{Ordre d'intégration (Théorique)}

\begin{enumerate}
  
\item On obtient ce résultat par récurrence. Par exemple, pour $k=2$, on a :

\begin{align*}
  (1-B)^2 (a_0+a_1t+a_2t^2) &= (1-B) ((a_0+a_1t+a_2t^2)\\ &\quad- (a_0+a_1(t-1)+a_2(t-1)^2)) \\
  &= (1-B) (a_1+a_2(2t+1)) \\
  &= (a_1+a_2(2t+1)) - (a_1+a_2(2(t-1)+1)) \\
  &= 2a_2
\end{align*}

En général, on obtient:

\begin{align*}
  (1-B)^k (a_0+a_1t+a_2t^2+\ldots+a_kt^k) &= (1-B)^{k-1} ((a_0+a_1t+a_2t^2+\ldots+a_kt^k)\\ &\quad- (a_0+a_1(t-1)+a_2(t-1)^2+\ldots+a_k(t-1)^k)) \\
  &= (1-B)^{k-1} (a_1 + 2a_2t+\ldots+a_k(t^k-(t-1)^k))
\end{align*}

On remarque qu'à chaque itération, le premier terme de la série disparait. Ainsi, après $k$ itérations, il ne restera que le terme en $a_k$ avec son coefficient, qui correspont à $k!$. On obtient ainsi la solution générale.

\item Une série est dite stationnaire lorsque chaque terme est un terme d'erreur dont la distribution est constante au fil du temps. Ainsi, la distribution de la différence de deux termes consécutifs de la série sera aussi constante au fil du temps. Par exemple:
  \begin{align*}
    \epsilon_t \sim N(0,\sigma^2)  \\
     \epsilon_t - \epsilon_{t-1} \sim N(0,2\sigma^2) \\
  \end{align*}

\end{enumerate}



\clearpage
\subsection{Inversion de processus d'ordre 1 (Théorique)}

\begin{enumerate}
\item Un processus AR(1) est défini par $y_t = \gamma_1y_{t-1} + \epsilon_t$. En développant le terme $y_{t-1}$, on obtient $y_t = \gamma_1^2y_{t-2}+\gamma_1\epsilon_{t-1}+\epsilon_t$. De manière récursive, on obtient $y_t = \gamma_1^{t}\epsilon_0 + \gamma_1^{t-1}\epsilon_1 + \ldots + \gamma_1\epsilon_{t-1} + \epsilon_t$. ainsi, en faisant tendre $t\to\infty$, on obtient une représentation MA($\infty$).
  
\item Un processus MA(1) est défini par $y_t = \epsilon_t - \theta_1\epsilon_{t-1}$. On cherche à substituer le terme $\epsilon_{t-1}$. On développe le terme précédent de la série: $y_{t-1} = \epsilon_{t-1} - \theta_1\epsilon_{t-2}$ et on substitue dans la première expression pour obtenir $y_t = \epsilon_t - \theta_1y_{t-1} - \theta_1^2\epsilon_{t-2}$. De manière récursive, on obtient $y_t = -\theta_1^ty_0-\theta_1^{t-1}y_1-\ldots-\theta_1y_{t-1}+\epsilon_{t}$. Lorsque $t\to\infty$, on obtient une représentation AR($\infty$).
\end{enumerate}


\clearpage
\subsection{Construction d'une série autorégressive (Calculatrice)}

On utilise la formule $y_t = \gamma_1y_{t-1} + \epsilon_t$.

\begin{center}
\begin{tabular}{|r|r|r|}
\hline
\multicolumn{1}{|l|}{} & \multicolumn{ 2}{c|}{$\gamma$} \\ \hline
\multicolumn{1}{|l|}{$N(0,1)$} & -0,5 & 0,5 \\ \hline
-1,21 & -1,2100 & -1,2100 \\ \hline
0,28 & 0,8850 & -0,3250 \\ \hline
1,08 & 0,6375 & 0,9175 \\ \hline
-2,35 & -2,6688 & -1,8913 \\ \hline
0,43 & 1,7644 & -0,5156 \\ \hline
0,51 & -0,3722 & 0,2522 \\ \hline
-0,57 & -0,3839 & -0,4439 \\ \hline
-0,55 & -0,3580 & -0,7720 \\ \hline
-0,56 & -0,3810 & -0,9460 \\ \hline
-0,89 & -0,6995 & -1,3630 \\ \hline
\end{tabular}
\end{center}

Les séries avec une corrélation négative ont tendance à aller dans la direction contraire des termes précédents alors que celles avec une corrélation positive ont tendance à aller dans la même direction que les termes précédents.

Le tableur \url{constructionserieAR.ods} contient les calculs effectués.

\subsection{Deux processus MA(2) (Théorique)}

Il suffit de calculer la fonction d'autorégression pour chaque processus MA(2).
Ensuite, on peut évaluer la fonction d'autocorrélation et comparer le résultat obtenu.

\begin{enumerate}
\item Premier processus avec:
  \begin{align*}
    \theta_1 = \theta_2 = \frac{1}{4}
  \end{align*}
  Fonction d'autocovariance
  \begin{align*}
    \gamma_0^{(1)} &= V[Y_t] \\
    &= V[e_t]+\frac{1}{16}V[e_{t-1}]+\frac{1}{16}V[e_{t-2}] \\
    &= (1+\frac{1}{8})\theta^2_e \\
    &= \frac{9}{8} \sigma^2_e \\
  \end{align*}
  \begin{align*}
    \gamma_1^{(1)} &= Cov(Y_t,Y_{t-1}) \\
    &= Cov(e_t - \frac{1}{4}e_{t-1} - \frac{1}{4}e_{t-2}, e_{t-1} - \frac{1}{4}e_{t-2} - \frac{1}{4}e_{t-3} )\\
    &= Cov(-\frac{1}{4}e_{t-1},e_{t-1}) + Cov(-\frac{1}{4}e_{t-2},-\frac{1}{4}e_{t-2}) \\
    &= (-\frac{1}{4}+(-\frac{1}{4})(-\frac{1}{4})) \sigma^2_e \\
    &= -\frac{3}{16} \sigma^2_e \\
  \end{align*}
  \begin{align*}
    \gamma_2^{(1)} &= Cov(Y_t,Y_{t-2}) \\
    &= Cov(e_t - \frac{1}{4}e_{t-1} - \frac{1}{4}e_{t-2}, e_{t-2} - \frac{1}{4}e_{t-3} - \frac{1}{4}e_{t-4} )\\
    &= Cov(-\frac{1}{4}e_{t-2},e_{t-2}) \\
    &= -\frac{1}{4} \sigma^2_e \\
  \end{align*}
  \begin{align*}
    \gamma_k^{(1)} &= 0, \qquad \forall k \geq 3
   \end{align*}
   Fonction d'autocorrélation
   \begin{align*}
     \rho_1^{(1)} &= \frac{\gamma_1^{(1)}}{\gamma_0^{(1)}} \\
     &= \frac{\frac{-3}{16}\sigma^2_e}{\frac{9}{8}\sigma^2_e} \\
     &= \frac{-1}{6} \\
   \end{align*}
   \begin{align*}
     \rho_2^{(1)} &= \frac{\gamma_2^{(1)}}{\gamma_0^{(1)}} \\
     &= \frac{\frac{-1}{4}\sigma^2_e}{\frac{9}{8}\sigma^2_e} \\
     &= \frac{-2}{9} \\
   \end{align*}
 \item Second processus avec:
   \begin{align*}
     \theta_1 = -1 \quad \theta_2 = 4
   \end{align*}
   Fonction d'autocovariance
   \begin{align*}
     \gamma_0^{(2)} &= V[Y_t] \\
     &= V[e_t]+V[e_{t-1}]+16V[e_{t-2}] \\
     &= 18 \sigma^2_e
   \end{align*}
   \begin{align*}
     \gamma_1^{(2)} &= Cov(Y_t,Y_{t-1}) \\
     &= Cov(e_t + e_{t-1} - 4e_{t-2}, e_{t-1} + e_{t-2} - 4e_{t-3}) \\
     &= Cov(e_{t-1},e_{t-1}) + Cov(-4e_{t-2},e_{t-2}) \\
     &= (1 + (-1)(4))\sigma^2_e \\
     &= -3\sigma^2_e \\
   \end{align*}
   \begin{align*}
     \gamma_2^{(2)} &= Cov(Y_t,Y_{t-2}) \\
     &= Cov(e_t + e_{t-1} - 4e_{t-2}, e_{t-2} + e_{t-3} - 4e_{t-4} )\\
     &= Cov(- 4e_{t-2},e_{t-2}) \\
     &= -4\sigma^2_e \\
   \end{align*}
   Fonction d'autocorrélation
   \begin{align*}
     \rho_1^{(2)} &= \frac{\gamma_1^{(2)}}{\gamma_0^{(2)}} \\
     &= \frac{-3\sigma^2_e}{18\sigma^2_e} \\
     &= \frac{-1}{6} \\
   \end{align*}
   \begin{align*}
     \rho_2^{(2)} &= \frac{\gamma_2^{(2)}}{\gamma_0^{(2)}} \\
     &= \frac{-4\sigma^2_e}{18\sigma^2_e} \\
     &= \frac{-2}{9} \\
   \end{align*}
 \end{enumerate}
 
 On remarque clairement que $\rho_1^{(1)} = \rho_1^{(2)}$ et $\rho_2^{(1)} = \rho_2^{(2)}$. 
 La fonction d'autocovariance vaut toujours 1 pour $\rho_1$ et vaut 0 ailleurs.
 
 \clearpage
 
 \includegraphics[height=7mm,keepaspectratio=true]{by-sa}\\%
Cette création est mise à disposition selon le contrat
\href{http://creativecommons.org/licenses/by-sa/2.5/ca/deed.fr}{%
  Paternité-Partage à l'identique 2.5 Canada} de Creative Commons
disponible à l'adresse \\
http://creativecommons.org/licenses/by-sa/2.5/ca/deed.fr \\

En vertu de ce contrat, vous êtes libre de :

\begin{itemize}
\item \textbf{partager} --- reproduire, distribuer et communiquer
  l'{\oe}uvre;
\item \textbf{remixer} --- adapter l'{\oe}uvre;
\item utiliser cette {\oe}uvre à des fins commerciales.
\end{itemize}

Selon les conditions suivantes:\\

  \begin{tabularx}{\linewidth}{@{}lX@{}}
    \raisebox{-9mm}[0mm][13mm]{%
      \includegraphics[height=11mm,keepaspectratio=true]{by}} &
    \textbf{Attribution} --- Vous devez attribuer l'{\oe}uvre de la
    manière indiquée par l'auteur de l'{\oe}uvre ou le titulaire des
    droits (mais pas d'une manière qui suggérerait qu'ils vous
    soutiennent ou
    approuvent votre utilisation de l'{\oe}uvre). \\
    \raisebox{-9mm}{\includegraphics[height=11mm,keepaspectratio=true]{sa}}
    & \textbf{Partage à l'identique} --- Si vous modifiez, transformez
    ou adaptez cette {\oe}uvre, vous n'avez le droit de distribuer
    votre création que sous une licence identique ou similaire à
    celle-ci.
  \end{tabularx}

%%% Local Variables: 
%%% mode: latex
%%% TeX-master: t
%%% End: 

\end{document}
